
\lussec 1. 引言

鉴于其重正化性质
[\cite{WilsonFlow},\cite{RenFlow}]，
并且它在格点规范理论中的应用在技术上十分直接，
杨--米尔斯梯度流使我们能以许多有趣的方式探究非阿贝尔
规范理论的动力学。
例如，该流可用于精确的标度设定，
而且它给出了拓扑（瞬子）扇区如何在
格点 QCD 的连续极限中显现的确切理解
[\cite{WilsonFlow}]\kern1pt\sfn{$\dagger$}{%
在格点规范理论中，当流方程右端与 Wilson 方格作用量的梯度一致时，
该流常被称为 ``Wilson 流''。
本文在连续理论与格点上均使用``梯度流''一词。}。
此外，正流时间下的观测量
是非微扰重正化与步长标度研究中
自然而然要考虑的量
[\cite{StepScaling}--\cite{FritzschRamos}]。

流方程中可以包含物质场，也可以不包含。
本文考虑将流向 QCD 夸克场的一个相当平凡的推广：
规范场的流方程保持不变，
而夸克场随流时间的演化
由一个规范协变的热方程决定（见第 2 节）。
夸克场的引入并不会使流的
理论分析与实际实现显著复杂化。特别地，
仿照文献~[\cite{RenFlow}]，可以证明：一旦 QCD 的参数按通常方式重正化，
正流时间下规范不变的局域场关联函数的重正化
只需要对随时间演化的夸克场作
一个乘法重正化。

作为示例，本文具体完成了推广流的两个应用：
其一是格点 QCD 中轴矢流重正化常数计算的
一种新策略，
其二是在很大程度上通过计算正流时间下
标量夸克密度的期望值来得到手征凝聚。
在这两个例子中，流的使用在技术上都被证明是有吸引力的。
例如，手征凝聚可以很容易地以高精度获得，
因为不需要加性重正化。

在接下来的两节中，我们在连续理论中讨论梯度流
向夸克场的推广。
由于流方程保持手征对称性，
随时间演化场的关联函数
满足简单的手征 Ward 恒等式。
如第 4 节所解释的，这些恒等式使人们能把
这些关联函数与轻赝标介子的物理联系起来。
随后是论文更偏技术的部分（第 5--7 节），
其中在格点 QCD 的框架下建立该流。特别地，
我们将证明与之相联系的 $4+1$ 维场论
容许一个良定义的局域格点正规化。
最后，通过在 $2+1$ 味 QCD 中的一个示例计算
（第 8、9 节）演示所提出的流的应用
在数值格点 QCD 中的可行性。
